\begin{tikzpicture}[
  >=Latex,
  manifold/.style={fill=blue!8, draw=blue!35!black, line width=0.75pt},
  gridline/.style={draw=blue!25!black, line width=0.35pt, opacity=0.55},
  geodesic/.style={draw=black!62, line width=0.85pt, dashed},
  meanconnect/.style={draw=black!62, line width=0.85pt, dashed},
  logvec/.style={-Latex, draw=orange!80!black, line width=0.85pt},
  meanvec/.style={-Latex, draw=purple!85!black, line width=1.25pt},
  expmap/.style={-Latex, draw=purple!85!black, line width=1.15pt, densely dashed},
  point/.style={circle, inner sep=1.7pt, fill=black},
  dest/.style={circle, inner sep=1.6pt, fill=red!75!black},
  bary/.style={circle, inner sep=2.0pt, fill=green!45!black},
  tang/.style={circle, inner sep=2.0pt, fill=purple!80!black},
  label/.style={font=\scriptsize, inner sep=1pt},  panel/.style={font=\bfseries\small}
]

% A reusable manifold patch shape, used both for drawing and clipping.
\def\ManifoldPatch{%
  (-3.15,-0.86) .. controls (-2.42,0.98) and (-0.72,1.48) .. (0.80,1.20)
  .. controls (2.48,0.90) and (3.38,-0.18) .. (2.78,-1.02)
  .. controls (1.70,-1.72) and (-1.92,-1.62) .. (-3.15,-0.86) -- cycle}

% ------------------------------------------------------------------
% Panel A: intrinsic barycentric projection
% ------------------------------------------------------------------
\begin{scope}[shift={(0,0)}]
  % manifold patch
  \draw[manifold] \ManifoldPatch;

  % Curves that should visually live on the manifold are clipped to the patch.
  \begin{scope}
    \clip \ManifoldPatch;
    % soft coordinate curves on the manifold
    \draw[gridline] (-2.65,-0.56) .. controls (-1.55,0.20) and (0.50,0.34) .. (2.55,-0.28);
    \draw[gridline] (-2.18,-1.04) .. controls (-0.90,-0.42) and (0.90,-0.38) .. (2.55,-0.80);
    \draw[gridline] (-2.02,0.67) .. controls (-0.55,0.98) and (1.13,0.77) .. (2.55,0.27);
    \draw[gridline] (-1.95,-1.18) .. controls (-1.66,-0.18) and (-1.63,0.55) .. (-1.33,1.10);
    \draw[gridline] (0.05,-1.34) .. controls (0.23,-0.35) and (0.23,0.50) .. (0.40,1.28);
    \draw[gridline] (1.92,-1.14) .. controls (1.58,-0.24) and (1.40,0.44) .. (1.19,1.02);

    % geodesics from intrinsic barycenter to destinations, clipped to stay inside M
    \coordinate (BAclip) at (0.78,0.02);
    \coordinate (yAoneclip) at (0.36,0.80);
    \coordinate (yAtwoclip) at (1.65,0.10);
    \coordinate (yAthreeclip) at (0.18,-0.74);
    \draw[geodesic] (BAclip) .. controls (0.60,0.30) and (0.42,0.55) .. (yAoneclip);
    \draw[geodesic] (BAclip) .. controls (1.12,0.15) and (1.38,0.11) .. (yAtwoclip);
    \draw[geodesic] (BAclip) .. controls (0.54,-0.28) and (0.32,-0.54) .. (yAthreeclip);
  \end{scope}

  % points
  \coordinate (xA) at (-2.28,-0.24);
  \coordinate (yAone) at (0.36,0.80);
  \coordinate (yAtwo) at (1.65,0.10);
  \coordinate (yAthree) at (0.18,-0.74);
  \coordinate (BA) at (0.78,0.02);

  \node[point] at (xA) {};
  \node[dest] at (yAone) {};
  \node[dest] at (yAtwo) {};
  \node[dest] at (yAthree) {};
  \node[bary] at (BA) {};

  \node[label, anchor=east] at ($(xA)+(-0.08,0.08)$) {$x$};
  \node[label, anchor=west] at ($(yAone)+(0.09,0.05)$) {$y_1$};
  \node[label, anchor=west] at ($(yAtwo)+(0.07,0.02)$) {$y_2$};
  \node[label, anchor=north] at ($(yAthree)+(0.14,-0.10)$) {$y_3$};
  \node[label, anchor=south west] at ($(BA)+(-0.02,-0.4)$) {$B_\pi(x)$};
\end{scope}

% ------------------------------------------------------------------
% Panel B: tangential barycentric projection
% ------------------------------------------------------------------
\begin{scope}[shift={(8.35,0)}]
  % manifold patch
  \draw[manifold] \ManifoldPatch;

  \begin{scope}
    \clip \ManifoldPatch;
    % coordinate curves
    \draw[gridline] (-2.65,-0.56) .. controls (-1.55,0.20) and (0.50,0.34) .. (2.55,-0.28);
    \draw[gridline] (-2.18,-1.04) .. controls (-0.90,-0.42) and (0.90,-0.38) .. (2.55,-0.80);
    \draw[gridline] (-2.02,0.67) .. controls (-0.55,0.98) and (1.13,0.77) .. (2.55,0.27);
    \draw[gridline] (-1.95,-1.18) .. controls (-1.66,-0.18) and (-1.63,0.55) .. (-1.33,1.10);
    \draw[gridline] (0.05,-1.34) .. controls (0.23,-0.35) and (0.23,0.50) .. (0.40,1.28);
    \draw[gridline] (1.92,-1.14) .. controls (1.58,-0.24) and (1.40,0.44) .. (1.19,1.02);

    % geodesic-like arcs from x to destinations, clipped to remain entirely on M
    \coordinate (xBclip) at (-2.28,-0.24);
    \coordinate (yBoneclip) at (0.36,0.80);
    \coordinate (yBtwoclip) at (1.65,0.10);
    \coordinate (yBthreeclip) at (0.18,-0.74);
    \draw[geodesic] (xBclip) .. controls (-1.55,0.50) and (-0.50,0.87) .. (yBoneclip);
    \draw[geodesic] (xBclip) .. controls (-0.95,0.05) and (0.35,0.17) .. (yBtwoclip);
    \draw[geodesic] (xBclip) .. controls (-1.10,-0.65) and (-0.32,-0.83) .. (yBthreeclip);
  \end{scope}

  \coordinate (xB) at (-2.28,-0.24);
  \coordinate (yBone) at (0.36,0.80);
  \coordinate (yBtwo) at (1.65,0.10);
  \coordinate (yBthree) at (0.18,-0.74);
  \coordinate (TB) at (0.32,0.08);

  % tangent plane: lower, visibly separated from the manifold, horizontally wide,
  % and vertically tall enough to contain the logarithmic displacement vectors.
  \coordinate (p0) at (-2.28,-2.28);
  \coordinate (p1) at (-3.74,-3.05);
  \coordinate (p2) at (2.40,-3.05);
  \coordinate (p3) at (2.84,-1.74);
  \coordinate (p4) at (-3.28,-1.74);
  \draw[fill=orange!10, draw=orange!55!black, line width=0.60pt, opacity=0.96]
    (p1) -- (p2) -- (p3) -- (p4) -- cycle;
  %\node[label, anchor=south east] at ($(p4)+(0.,0.)$) {$T_xM$};
  \draw[draw=black!35, densely dotted] (xB) -- (p0);
  \node[label, anchor=east, black!65] at ($(p0)+(-0.02,0.00)$) {$0_x$};

  % destinations and tangential projection on manifold
  \node[point] at (xB) {};
  \node[dest] at (yBone) {};
  \node[dest] at (yBtwo) {};
  \node[dest] at (yBthree) {};
  \node[tang] at (TB) {};

  \node[label, anchor=east] at ($(xB)+(-0.08,0.08)$) {$x$};
  \node[label, anchor=west] at ($(yBone)+(0.09,0.05)$) {$y_1$};
  \node[label, anchor=west] at ($(yBtwo)+(0.07,0.02)$) {$y_2$};
  \node[label, anchor=north] at ($(yBthree)+(0.14,-0.10)$) {$y_3$};
  \node[label, anchor=south west] at ($(TB)+(0.07,0.05)$) {$\widetilde B_\pi(x)$};

  % log vectors in tangent plane
  \coordinate (vone) at (-0.78,-2.08);
  \coordinate (vtwo) at (0.60,-2.34);
  \coordinate (vthree) at (-0.98,-2.56);
  \coordinate (vmean) at (-0.09,-2.21);

  \draw[logvec] (p0) -- (vone);
  \draw[logvec] (p0) -- (vtwo);
  \draw[logvec] (p0) -- (vthree);
  \draw[meanvec] (p0) -- (vmean);

  % dashed connections from the log-vector endpoints to their average
  \draw[meanconnect] (vone) -- (vmean);
  \draw[meanconnect] (vtwo) -- (vmean);
  \draw[meanconnect] (vthree) -- (vmean);

  \node[label, anchor=south east] at ($(vone)+(-0.2,0)$) {$\log_x(y_1)$};
  \node[label, anchor=west] at ($(vtwo)+(0.06,0.00)$) {$\log_x(y_2)$};
  \node[label, anchor=north east] at ($(vthree)+(-0.04,-0.03)$) {$\log_x(y_3)$};
  \node[label, anchor=south] at ($(vmean)+(-0.1,0.10)$) {$m_\pi(x)$};

  % exponential map from averaged vector back to manifold. The curve is routed
  % to the right so that it remains visually separated from y_3.
  \draw[expmap] (vmean)
    .. controls (1.08,-1.58) and (1.30,-0.50) ..
    node[label, pos=0.54, anchor=west, xshift=0.14cm, yshift=0.01cm, purple!85!black] {$\exp_x$}
    (TB);
\end{scope}

\end{tikzpicture}